首先，衷心感謝指導教授劉怡芬老師在本研究全程中給予的悉心引領。感謝老師給予我踏入語音處理與深度學習研究領域的機會，使我得以在這個充滿挑戰與魅力的方向上持續深耕。老師慷慨撥出大量個人研究時間，耐心聆聽每一次的報告與進度說明，並在每個岔路口給予明確的方向指引；對於我諸多天馬行空的想法，老師亦願意逐一與我深入討論，使研究得以在反覆推敲中逐步聚焦。從每一次的開會與討論中，我從老師身上看到了對於學術研究應有的態度與熱忱；在老師百忙之中仍抽空與我討論、指引方向的過程裡，我不僅在專業知識上獲得了長足的成長，更深刻體會到做研究所需的嚴謹與堅持，這些寶貴的經驗，將是我往後繼續深耕學術道路上最重要的養分。

其次，感謝研究夥伴陳柄佑在研究過程中的一路同行。研究夥伴全程出席每次討論會議，並在即時深偽偵測系統的建構上投入大量心力——涵蓋以 Next.js 實作的前端互動介面、以 FastAPI 搭建的後端推論服務，乃至以 Docker 容器化整套部署環境，使系統得以在複雜的深度學習依賴環境下穩定運行。在資料集方面，研究夥伴協助蒐集並整理了本研究所使用的多組語料，包含 TTS 相關技術資料的查閱彙整，以及各訓練與測試資料集的下載、格式統一與前處理，為後續實驗的順利推進奠定了基礎。在研究最為艱辛的時刻，也始終願意傾聽，是本研究得以完成不可或缺的夥伴。

\begin{flushright}
\mbox{學生 王兆偉}
\end{flushright}
